\documentclass{amsart}
% For dealing with references we use the comment environment
\usepackage{verbatim}
\newenvironment{reference}{\comment}{\endcomment}
%\newenvironment{reference}{}{}
\newenvironment{slogan}{\comment}{\endcomment}
\newenvironment{history}{\comment}{\endcomment}

% For commutative diagrams we use Xy-pic
\usepackage[all]{xy}

% We use 2cell for 2-commutative diagrams.
\xyoption{2cell}
\UseAllTwocells

% We use multicol for the list of chapters between chapters
\usepackage{multicol}

% This is generally recommended for better output
\usepackage{lmodern}
\usepackage[T1]{fontenc}

% For cross-file-references

% Package for hypertext links:
\usepackage{hyperref}

% For any local file, say "hello.tex" you want to link to please

% Theorem environments.
%
\theoremstyle{plain}
\newtheorem{theorem}[subsection]{Theorem}
\newtheorem{proposition}[subsection]{Proposition}
\newtheorem{lemma}[subsection]{Lemma}

\theoremstyle{definition}
\newtheorem{definition}[subsection]{Definition}
\newtheorem{example}[subsection]{Example}
\newtheorem{exercise}[subsection]{Exercise}
\newtheorem{situation}[subsection]{Situation}

\theoremstyle{remark}
\newtheorem{remark}[subsection]{Remark}
\newtheorem{remarks}[subsection]{Remarks}

\numberwithin{equation}{subsection}

% Macros
%
\def\lim{\mathop{\mathrm{lim}}\nolimits}
\def\colim{\mathop{\mathrm{colim}}\nolimits}
\def\Spec{\mathop{\mathrm{Spec}}}
\def\Hom{\mathop{\mathrm{Hom}}\nolimits}
\def\Ext{\mathop{\mathrm{Ext}}\nolimits}
\def\SheafHom{\mathop{\mathcal{H}\!\mathit{om}}\nolimits}
\def\SheafExt{\mathop{\mathcal{E}\!\mathit{xt}}\nolimits}
\def\Sch{\mathit{Sch}}
\def\Mor{\mathop{\mathrm{Mor}}\nolimits}
\def\Ob{\mathop{\mathrm{Ob}}\nolimits}
\def\Sh{\mathop{\mathit{Sh}}\nolimits}
\def\NL{\mathop{N\!L}\nolimits}
\def\CH{\mathop{\mathrm{CH}}\nolimits}
\def\proetale{{pro\text{-}\acute{e}tale}}
\def\etale{{\acute{e}tale}}
\def\QCoh{\mathit{QCoh}}
\def\Ker{\mathop{\mathrm{Ker}}}
\def\Im{\mathop{\mathrm{Im}}}
\def\Coker{\mathop{\mathrm{Coker}}}
\def\Coim{\mathop{\mathrm{Coim}}}

% Boxtimes
%
\DeclareMathSymbol{\boxtimes}{\mathbin}{AMSa}{"02}

%
% Macros for moduli stacks/spaces
%
\def\QCohstack{\mathcal{QC}\!\mathit{oh}}
\def\Cohstack{\mathcal{C}\!\mathit{oh}}
\def\Spacesstack{\mathcal{S}\!\mathit{paces}}
\def\Quotfunctor{\mathrm{Quot}}
\def\Hilbfunctor{\mathrm{Hilb}}
\def\Curvesstack{\mathcal{C}\!\mathit{urves}}
\def\Polarizedstack{\mathcal{P}\!\mathit{olarized}}
\def\Complexesstack{\mathcal{C}\!\mathit{omplexes}}
% \Pic is the operator that assigns to X its picard group, usage \Pic(X)
% \Picardstack_{X/B} denotes the Picard stack of X over B
% \Picardfunctor_{X/B} denotes the Picard functor of X over B
\def\Pic{\mathop{\mathrm{Pic}}\nolimits}
\def\Picardstack{\mathcal{P}\!\mathit{ic}}
\def\Picardfunctor{\mathrm{Pic}}
\def\Deformationcategory{\mathcal{D}\!\mathit{ef}}

\setlength{\emergencystretch}{3em}
\begin{document}
\title[Stacks Project, Tag 0BK5]{711 --- Spectral convergence under filtered-cohomology vanishing and stabilization}
\maketitle
\noindent Copyright (C) 2005--2025 Johan de Jong. Stacks Project authors. Source: \href{https://stacks.math.columbia.edu/tag/0BK5}{Tag 0BK5}.

\noindent Editorial difficulty note (not author text). Difficulty H2: The author proves boundedness and identifies limiting cycles and boundaries directly using cohomological vanishing at one end and stabilization at the other. Controlling weak convergence and abutment separately for a filtered complex is above routine spectral-sequence exercises..

\noindent Editorial provenance note (not author text). History: excerpt from revision \texttt{a0444\allowbreak 6e57e\allowbreak c1fbc\allowbreak 252a8\allowbreak 71afc\allowbreak ec775\allowbreak 2fb28\allowbreak 07b14}, homology.tex, lines 6463--6542. Reference presentation was adapted; original-author context and core support blocks were retained or added. The mathematical wording is unchanged. Licensed under GNU FDL 1.2 or later; no invariant sections or cover texts. See ../licenses/COPYING. Context blocks reproduce author definitions and supporting results; they are not additional counted problems.

\begin{definition}
\label{definition-filtration-cohomology-filtered-complex}
Let $\mathcal{A}$ be an abelian category.
Let $(K^\bullet, F)$ be a filtered complex of $\mathcal{A}$.
The {\it induced filtration} on $H^n(K^\bullet)$ is the filtration defined
by $F^pH^n(K^\bullet) = \Im(H^n(F^pK^\bullet) \to H^n(K^\bullet))$.
\end{definition}

\begin{definition}
\label{definition-limit-spectral-sequence}
Let $\mathcal{A}$ be an abelian category.
Let $(E_r, d_r)_{r \geq 1}$ be a spectral sequence.
\begin{enumerate}
\item If the subobjects $Z_{\infty} = \bigcap Z_r$
and $B_{\infty} = \bigcup B_r$ of $E_1$ exist then we define
the {\it limit}\footnote{This notation is not universally accepted. In some
references an additional pair of subobjects
$Z_\infty$ and $B_\infty$ of $E_1$ such that
$0 = B_1 \subset B_2 \subset \ldots \subset B_\infty \subset Z_\infty
\subset \ldots \subset Z_2 \subset Z_1 = E_1$
is part of the data comprising a spectral sequence!}
of the spectral sequence to be the object
$E_{\infty} = Z_{\infty}/B_{\infty}$.
\item We say that the spectral sequence {\it degenerates at $E_r$}
if the differentials $d_r, d_{r + 1}, \ldots$ are all zero.
\end{enumerate}
\end{definition}

\begin{definition}
\label{definition-bounded-ss}
Let $\mathcal{A}$ be an abelian category. Let $(E_r, d_r)_{r \geq r_0}$
be a spectral sequence of bigraded objects of
$\mathcal{A}$ with $d_r$ of bidegree $(r, -r + 1)$.
We say such a spectral sequence is
\begin{enumerate}
\item {\it regular} if for all $p, q \in \mathbf{Z}$ there is
a $b = b(p, q)$ such that the maps
$d_r^{p, q} : E_r^{p, q} \to E_r^{p + r, q - r + 1}$ are zero for $r \geq b$,
\item {\it coregular} if for all $p, q \in \mathbf{Z}$ there is a
$b = b(p, q)$ such that the maps
$d_r^{p - r, q + r - 1} : E_r^{p - r, q + r - 1} \to E_r^{p, q}$
are zero for $r \geq b$,
\item {\it bounded} if for all $n$
there are only a finite number of nonzero $E_{r_0}^{p, n - p}$,
\item {\it bounded below} if for all $n$ there is a $b = b(n)$ such that
$E_{r_0}^{p, n - p} = 0$ for $p \geq b$.
\item {\it bounded above} if for all $n$ there is a $b = b(n)$ such that
$E_{r_0}^{p, n - p} = 0$ for $p \leq b$.
\end{enumerate}
\end{definition}

\begin{lemma}
\label{lemma-compute-cohomology-filtered-complex}
Let $\mathcal{A}$ be an abelian category. Let $(K^\bullet, F)$ be a filtered
complex of $\mathcal{A}$. If $Z_\infty^{p, q}$ and $B_\infty^{p, q}$ exist
(see proof), then
\begin{enumerate}
\item the limit $E_\infty$ exists and is a bigraded object having
$E_\infty^{p, q} = Z_\infty^{p, q}/B_\infty^{p, q}$ in bidegree $(p, q)$,
\item the $p$th graded part $\text{gr}^pH^n(K^\bullet)$ of the
$n$th cohomology object of $K^\bullet$ is a subquotient of
$E_\infty^{p, n - p}$.
\end{enumerate}
\end{lemma}

\begin{proof}
The objects $Z_\infty$, $B_\infty$, and the limit
$E_\infty = Z_\infty/B_\infty$ of
Definition \ref{definition-limit-spectral-sequence}
are bigraded objects of $\mathcal{A}$ by our construction of the
spectral sequence in Lemma \href{https://stacks.math.columbia.edu/tag/012M}{lemma spectral sequence filtered complex (Tag 012M)}.
Since $Z_r^{p, q}$, resp.\ $B_r^{p, q}$ is the bidegree $(p, q)$
part of $Z_r$, resp.\ $B_r$, if we assume that
$$
Z_\infty^{p, q} = \bigcap\nolimits_r Z_r^{p, q} =
\bigcap\nolimits_r
\frac{F^pK^{p + q} \cap d^{-1}(F^{p + r}K^{p + q + 1}) + F^{p + 1}K^{p + q}}
{F^{p + 1}K^{p + q}}
$$
and
$$
B_\infty^{p, q} = \bigcup\nolimits_r B_r^{p, q} =
\bigcup\nolimits_r
\frac{F^pK^{p + q} \cap d(F^{p - r + 1}K^{p + q - 1}) + F^{p + 1}K^{p + q}}
{F^{p + 1}K^{p + q}}
$$
exist, then $Z_\infty$ and $B_\infty$ exist with bidegree $(p, q)$
parts $Z_\infty^{p, q}$ and $B_\infty^{p, q}$ (follows from an elementary
argument about unions and intersections of bigraded objects). Thus
$$
E_\infty^{p, q} =
\frac{\bigcap_r (F^pK^{p + q} \cap d^{-1}(F^{p + r}K^{p + q + 1})
+ F^{p + 1}K^{p + q})}
{\bigcup_r (F^pK^{p + q} \cap d(F^{p - r + 1}K^{p + q - 1})
+ F^{p + 1}K^{p + q})}.
$$
where the top and the bottom exist. With $n = p + q$ we have
\begin{equation}
\label{equation-on-top-bigraded}
\Ker(d) \cap F^pK^{n} + F^{p + 1}K^{n}
\subset
\bigcap\nolimits_r
\left(
F^pK^{n} \cap d^{-1}(F^{p + r}K^{n + 1}) + F^{p + 1}K^{n}
\right)
\end{equation}
and
\begin{equation}
\label{equation-at-bottom-bigraded}
\bigcup\nolimits_r
\left(
F^pK^{n} \cap d(F^{p - r + 1}K^{n - 1}) + F^{p + 1}K^{n}
\right)
\subset
\Im(d) \cap F^pK^{n} + F^{p + 1}K^{n}.
\end{equation}
Thus a subquotient of $E_\infty^{p, q}$ is
$$
\frac{\Ker(d) \cap F^pK^{n} + F^{p + 1}K^n}
{\Im(d) \cap F^pK^{n} + F^{p + 1}K^{n}} =
\frac{\Ker(d) \cap F^pK^n}{\Im(d) \cap F^pK^n + \Ker(d) \cap F^{p + 1}K^n}
$$
Comparing with (\href{https://stacks.math.columbia.edu/tag/0BDT}{equation graded cohomology (Tag 0BDT)}) we conclude.
\end{proof}


\begin{lemma}
\label{lemma-filtered-complex-ss-converges}
Let $\mathcal{A}$ be an abelian category. Let $(K^\bullet, F)$ be a filtered
complex of $\mathcal{A}$. The associated spectral sequence
\begin{enumerate}
\item weakly converges to $H^*(K^\bullet)$ if and only if for every
$p, q \in \mathbf{Z}$ we have equality in equations
(\ref{equation-at-bottom-bigraded}) and (\ref{equation-on-top-bigraded}),
\item abuts to $H^*(K)$ if and only if it weakly converges to $H^*(K^\bullet)$
and we have
$\bigcap_p (\Ker(d) \cap F^pK^n + \Im(d) \cap K^n) = \Im(d) \cap K^n$
and
$\bigcup_p (\Ker(d) \cap F^pK^n + \Im(d) \cap K^n) = \Ker(d) \cap K^n$.
\end{enumerate}
\end{lemma}

\begin{proof}
Immediate from the discussions above.
\end{proof}


\begin{lemma}
\label{lemma-ss-converges-trivial}
Let $\mathcal{A}$ be an abelian category. Let $(K^\bullet, F)$ be a
filtered complex of $\mathcal{A}$. Assume
\begin{enumerate}
\item for every $n$ there exist $p_0(n)$ such that
$H^n(F^pK^\bullet) = 0$ for $p \geq p_0(n)$,
\item for every $n$ there exist $p_1(n)$ such that
$H^n(F^pK^\bullet) \to H^n(K^\bullet)$ is an isomorphism
for $p \leq p_1(n)$.
\end{enumerate}
Then
\begin{enumerate}
\item the spectral sequence associated to $(K^\bullet, F)$ is bounded,
\item the filtration on each $H^n(K^\bullet)$ is finite,
\item the spectral sequence associated to $(K^\bullet, F)$ converges
to $H^*(K^\bullet)$.
\end{enumerate}
\end{lemma}

\begin{proof}
Fix $n$. Using the long exact cohomology sequence associated to
the short exact sequence of complexes
$$
0 \to F^{p + 1}K^\bullet \to F^pK^\bullet \to \text{gr}^pK^\bullet \to 0
$$
we find that $E_1^{p, n - p} = 0$ for $p \geq \max(p_0(n), p_0(n + 1))$ and
$p < \min(p_1(n), p_1(n + 1))$. Hence the spectral sequence is bounded
(Definition \ref{definition-bounded-ss}). This proves (1).

\medskip\noindent
It is clear from the assumptions and
Definition \ref{definition-filtration-cohomology-filtered-complex}
that the filtration on $H^n(K^\bullet)$ is finite. This proves (2).

\medskip\noindent
Next we prove that the spectral sequence weakly converges to
$H^*(K^\bullet)$ using
Lemma \ref{lemma-filtered-complex-ss-converges}.
Let us show that we have equality in (\ref{equation-on-top-bigraded}).
Namely, for $p + r > p_0(n + 1)$ the map
$$
d : F^pK^{n} \cap d^{-1}(F^{p + r}K^{n + 1}) \to F^{p + r}K^{n + 1}
$$
ends up in the image of $d : F^{p + r}K^n \to F^{p + r}K^{n + 1}$
because the complex $F^{p + r}K^\bullet$ is exact in degree $n + 1$.
We conclude that $F^pK^{n} \cap d^{-1}(F^{p + r}K^{n + 1}) =
d(F^{p + r}K^n) + \Ker(d) \cap F^pK^n$. Hence for such $r$ we have
$$
\Ker(d) \cap F^pK^{n} + F^{p + 1}K^{n} =
F^pK^{n} \cap d^{-1}(F^{p + r}K^{n + 1}) + F^{p + 1}K^{n}
$$
which proves the desired equality. To show that we have equality in
(\ref{equation-at-bottom-bigraded}) we use that for $p - r + 1 < p_1(n - 1)$
we have
$$
d(F^{p - r + 1}K^{n - 1}) = \Im(d) \cap F^{p - r + 1}K^n
$$
because the map $F^{p - r + 1}K^\bullet \to K^\bullet$ induces an
isomorphism on cohomology in degree $n - 1$. This shows that
we have
$$
F^pK^{n} \cap d(F^{p - r + 1}K^{n - 1}) + F^{p + 1}K^{n} =
\Im(d) \cap F^pK^{n} + F^{p + 1}K^{n}
$$
for such $r$ which proves the desired equality.

\medskip\noindent
To see that the spectral sequence abuts to $H^*(K^\bullet)$ using
Lemma \ref{lemma-filtered-complex-ss-converges} we have to show that
$\bigcap_p (\Ker(d) \cap F^pK^n + \Im(d) \cap K^n) = \Im(d) \cap K^n$
and
$\bigcup_p (\Ker(d) \cap F^pK^n + \Im(d) \cap K^n) = \Ker(d) \cap K^n$.
For $p \geq p_0(n)$ we have
$\Ker(d) \cap F^pK^n + \Im(d) \cap K^n = \Im(d) \cap K^n$
and for $p \leq p_1(n)$ we have
$\Ker(d) \cap F^pK^n + \Im(d) \cap K^n = \Ker(d) \cap K^n$.
Combining weak convergence, abutment, and boundedness we see
that (2) and (3) are true.
\end{proof}
\end{document}
